\textit{Proof.}
Put \(v=\sigma^2/n\).  For a Gaussian shift with mean \(\mu_g\) and covariance
\(vI_{12}\), its likelihood ratio relative to \(P_0\) is
\[
L_g(y)=\exp\!\left\{v^{-1}\mu_g^{\mathsf T}y
-\frac{1}{2v}\lVert\mu_g\rVert_2^2\right\}.
\]
Completing the square under \(Y\sim P_0\) gives, for every \(g,g'\),
\[
E_{P_0}[L_g(Y)L_{g'}(Y)]
=\exp\!\left(v^{-1}\langle\mu_g,\mu_{g'}\rangle\right).
\]
Since \(d\overline P_{h,q}/dP_0=64^{-1}\sum_gL_g\), expansion of the square
proves
\[
1+\chi^2(\overline P_{h,q},P_0)
=64^{-2}\sum_{g,g'\in\mathcal G}
\exp\!\left(\frac n{\sigma^2}
\langle\mu_g(h,q),\mu_{g'}(h,q)\rangle\right).
\]

For \(h=q=ta\), write \(\mu_g(ta,ta)=ta\nu_g\).  The orbit-membership
calculation gives \(64^{-1}\sum_g\nu_g=0\).  Therefore the coefficient of
\(t^2\) in the preceding exponential sum is
\[
\frac{na^2}{\sigma^2},64^{-2}
\sum_{g,g'}\langle\nu_g,\nu_{g'}\rangle
=\frac{na^2}{\sigma^2}
\left\lVert64^{-1}\sum_g\nu_g\right\rVert_2^2=0.
\]
For each diagonal term, the three observed blocks have squared Frobenius norms
\(2t^2a^2\), \(2t^2a^2\), and \(4t^2a^2\), respectively.  Thus
\(\lVert\mu_g(ta,ta)\rVert_2^2=8t^2a^2\).  Keeping only the sixty-four
diagonal, nonnegative terms in the exact Gram sum yields
\[
1+\chi^2(\overline P_{ta,ta},P_0)
\geq64^{-1}\exp(8na^2t^2/\sigma^2).
\]
If \(ta\asymp n^{-1/4}\), the exponent is of order \(n^{1/2}\), and the
right side diverges.  Hence cancellation of the quadratic Taylor coefficient
does not give bounded chi-square affinity at that scaling. \(\square\)